% ---
% id: meat-core/ass:fixed
% kind: ass
% title: "Protocol parameters fixed before the run"
% summary: >-
%   The acceptance set, the round count, the key length, the transcript
%   length, the error budgets, the authentication parameters and the Rényi
%   order are fixed before the run and applied unchanged. The key length is
%   that of a fixed-length protocol, not one chosen from the run's data.
% prominence: standard
% status: accepted
% uses: [def:postprocessing-model, def:authenticated-channel]
% assumes: []
% refs:
%   - "tupkary2025survey, Sec. 2.2"
% provenance: newly-derived
% short_id: A-FIXED
% verification: never
% category: protocol-structure
% affects: [secrecy, correctness]
% user_statement: >-
%   Every parameter the key length rests on is fixed before the run starts:
%   the acceptance test, the number of rounds, the key length, the
%   error-correction budget λ_EC, the error budgets ε_PA and ε_EV, and the
%   authentication parameters ε_auth and k_auth. The Rényi order α at which
%   the key length is computed is fixed before the run too; choosing it by
%   comparing the key lengths your declared model gives at several orders,
%   as the run tier's outer search over α does, reads no data of the run and is such
%   a choice. None of them may be chosen or adjusted after looking at the
%   run's data, for example by relaxing the acceptance threshold or picking
%   the key length from the observed error rate. That would be a
%   variable-length protocol, which this analysis does not cover.
% lean:
% ---

A fixed-length protocol outputs a key of a length determined before the
protocol begins, or aborts \cite[Sec.~2.2]{tupkary2025survey}; choosing the
length or the test from the observed data makes it a variable-length
protocol, which needs a different analysis. The R\'enyi order $\alpha$ is a
parameter of the proof, not of the protocol: the run never reads it, and it
enters only the computation of $\ell$. Clause (iii) fixes it before the run
all the same, so that the order, and with it the length computed at it, is
never a function of the run's data; an order chosen by comparing the
lengths that the \tref{declared model} gives at several orders satisfies it.
No proof consumes clause (iii): the key-length theorem holds at each
$\alpha\in(1,2)$, and the clause is stated so that the attestation names
every input of $\ell$.

\begin{assumption}\label{ass:fixed}
(i) The acceptance set $\acSet$, the round count $n$, the output length
$\ell$, the transcript length $\lEC$, the error budgets $\epa$ and $\eev$ of
Definition~\ref{def:postprocessing-model}, and the authentication parameters
$\varepsilon_{\mathrm{auth}}$ and $k_{\mathrm{auth}}$, the declared
authentication error bound and key consumption
(Definition~\ref{def:authenticated-channel}\,(b)), are fixed before the
first round of the run;
(ii) the run applies them unchanged: none of them is chosen, adjusted or
recomputed from data observed in the run;
(iii) the R\'enyi order $\alpha\in(1,2)$ at which $\ell$ is computed is
fixed before the first round of the run, from data available before it, and
is not chosen, adjusted or recomputed from data observed in the run.
\end{assumption}
